% =============================================================================
% cap3_pipeline/05_fase4_reid.tex
% Seção 3.5 -- Fase IV: Reidentificação individual de Felis catus
%   Modelos: PetFace (face) + método de re-ID por partes do corpo
%   (Akbar 2025) (flanco/padrão de pelagem/ear-tipping).
%   Função: dado um quadro recortado e classificado como Felis catus,
%   atribuir um id individual contra o catálogo da colônia, ou marcar
%   como "novo indivíduo" para curadoria humana.
%   Texto corrido, impessoal, sem listas no corpo.
% =============================================================================
\section{Fase IV --- Reidentificação individual}
\label{sec:pipeline-fase4}

A Fase~IV é a tarefa-fim do pipeline e a justificativa de existência de todo o sistema. Recebida do estágio anterior uma coleção depurada de quadros confiantemente classificados como \estr{Felis catus}, seu objetivo é, para cada quadro, atribuir um identificador individual ao animal --- ou registrar, com a devida cautela, que o indivíduo observado não é compatível com nenhum dos cadastrados, hipótese que aciona o protocolo de curadoria humana descrito ao final desta seção. A reidentificação individual é o passo que transforma um detector genérico de fauna em um sistema de monitoramento populacional, permitindo responder às perguntas que efetivamente interessam ao manejo da colônia: quantos indivíduos distintos visitam cada ponto, com que frequência, com que sobreposição entre pontos e com que continuidade temporal.

\subsection*{Decomposição da tarefa: face e flanco}

A literatura recente em re-ID de carnívoros pequenos converge em recomendar que a tarefa não seja tratada como problema único de \emph{embedding} corporal global, mas decomposta em conjuntos de \emph{features} discriminativas distintas, combináveis em decisão final ponderada. Para \estr{Felis catus} em ambiente comunitário, dois conjuntos têm sustentação empírica direta: a face, com sua estrutura ocular, nasal e auricular específica, e o flanco, com seu padrão de pelagem, contorno corporal e --- no caso de colônias submetidas a CED --- a marcação por corte reto da ponta da orelha (\emph{ear-tipping}). Essa decomposição é a abordagem proposta por \citeonline{akbar2025bodypart} para re-ID por partes do corpo em felinos pequenos, alinhada à disponibilização do dataset \modeloPetFace, focado em reconhecimento facial de animais de companhia~\cite{shinoda2024petface}. A robustez da pelagem como atributo de re-ID é validada em larga escala no dataset \modeloWildlife~\cite{adam2024wildlifereid10k}.

Em decorrência, a Fase~IV é composta por dois subprocessos paralelos sobre o mesmo quadro recortado: re-ID por face, baseado em \modeloPetFace, e re-ID por flanco, baseado no método de partes do corpo de \citeonline{akbar2025bodypart}. Cada subprocesso produz um ranqueamento dos indivíduos do catálogo da colônia ordenados por similaridade ao quadro analisado; a decisão final combina os dois ranqueamentos por regra simples de fusão (média ponderada de \emph{scores} normalizados, ou eleição quando há concordância no topo do \RankK{1}), priorizando o subprocesso cuja entrada esteja em condições geometricamente mais favoráveis (face frontal disponível versus apenas flanco visível).

\subsection*{Modelos pré-treinados e catálogo da colônia}

Em coerência com a decisão de escopo da Seção~\ref{sec:pipeline-visao-geral}, ambos os subprocessos são executados com modelos pré-treinados sem \emph{fine-tuning} sobre dados do Campus~2. A \emph{pipeline} pública do \modeloPetFace~e a implementação de referência de \citeonline{akbar2025bodypart} são aplicadas \emph{as-is} sobre os recortes da Fase~III, produzindo \emph{embeddings} (vetores de características densos) que são comparados, por distância no espaço de embedding, contra o catálogo de embeddings dos indivíduos cadastrados. O catálogo é construído a partir de fotografias cadastrais dos animais já submetidos ao protocolo CED da AEX Gatosdoc2 (resumido na Seção~\ref{sec:caracterizacao-pontos}), processadas uma única vez pelos mesmos modelos e armazenadas como referência permanente no servidor. A atualização do catálogo --- inclusão de novos indivíduos, retirada de animais que deixaram a colônia --- é tarefa operacional dos voluntários, mediada pela mesma interface da fila de validação humana da Fase~III. Eventual \emph{fine-tuning} leve dos dois modelos sobre o conjunto cadastral local fica registrado como trabalho futuro no Capítulo~\ref{cap:conclusao}.

\subsection*{Saída da fase: \RankK{} e tratamento de novos indivíduos}

Para cada quadro processado, a Fase~IV produz uma lista ranqueada dos \(k\) indivíduos do catálogo mais similares ao animal observado, acompanhada de seus \emph{scores} de similaridade fundidos entre face e flanco. A política de decisão é cautelosa, em linha com a postura conservadora das fases anteriores. Quando o indivíduo do \RankK{1} apresenta similaridade acima de um limiar configurável e margem confortável sobre o \RankK{2}, o pipeline atribui automaticamente o identificador correspondente e o adiciona ao histórico do indivíduo. Quando a similaridade do \RankK{1} é alta mas a margem para o \RankK{2} é estreita --- caso típico em colônias com forte parentesco e padrões de pelagem similares ---, o registro é encaminhado para curadoria humana, com a lista ranqueada exibida ao voluntário para decisão. Quando nenhum indivíduo do catálogo supera o limiar mínimo, o registro é marcado como ``novo indivíduo candidato'' e fica retido para curadoria, em que a equipe da AEX Gatosdoc2 decide se corresponde a um animal recém-incorporado ainda não cadastrado ou a um falso negativo do filtro multi-espécie da Fase~III.

O desenho cumpre dois objetivos. O primeiro é evitar que o pipeline invente identidades: a atribuição automática só ocorre em alta confiança, e os demais casos seguem para curadoria humana. O segundo é manter o catálogo auditável: cada identidade automaticamente atribuída fica acompanhada do \emph{score} de similaridade e da margem para o segundo colocado, e cada decisão humana fica registrada com data e identificador do voluntário, formando uma trilha mínima de auditoria que será incorporada à implementação futura do sistema. A confiabilidade dos modelos isoladamente e do pipeline integrado é estimada no Capítulo~\ref{cap:metodologia} por meio de \mAP, \Fscore~e \RankK{k} apurados sobre datasets públicos representativos das condições esperadas no Campus~2.
